%% kapa: exam=1 round=20 year=2009 period=2 form=A subject=appraisal_law subject_ko=부동산관계법규 seq=1 number=41
%% source: Q-Net 감정평가사 기출문제 (cst00309) artlSeq=5024178 file=extracted/1008/5024178/2009년 제20회 감정평가사 1차시험 기출문제/2교시.hwp
%% conversion: kordoc 4.24.51 -> tools/gen1.py (auto); review=proofread-sonnet; flags=none
\begin{question}{41}
국토의 계획 및 이용에 관한 법령상 토지거래계약허가를 거부하여야 하는 경우로 옳지 \uline{않은} 것은?
\begin{choices}
  \item 토지이용목적이 자기의 거주용 주택용지로 이용하고자 하는 것이 아닌 경우
  \item 토지이용목적이 공익사업을 위한 토지 등의 취득 및 보상에 관한 법률상 토지를 수용 또는 사용을 할 수 있는 사업시행자가 그 사업을 시행하기 위하여 필요한 것이 아닌 경우
  \item 토지거래계약을 체결하고자 하는 자의 토지이용목적이 도시계획이나 그 밖에 토지의 이용 및 관리에 관한 계획에 맞지 아니한 경우
  \item 토지의 이용목적으로 보아 그 면적이 적합하지 아니하다고 인정되는 경우
  \item 허가구역 지정 당시 그 구역에서 사업을 시행하고 있는 자가 그 사업에 이용하려는 것인 경우
\end{choices}
\end{question}
\answer{5}
